\section{F12 · Minste kvadraters metode}

\textbf{Kjernen:} Flere målinger enn parametre gir et overbestemt system som ikke kan løses eksakt. Vi minimerer summen av kvadrerte avvik. Er modellen lineær i parametrene, gir det normallikningene (eller QR, eller $A^\dagger$). Er modellen ikke-lineær, bruker vi Gauss--Newton eller Levenberg--Marquardt.

\subsection{Lineære modeller}
En modell $y=m(t;\vb x)$ der parametrene $\vb x\in\R^n$ inngår \emph{lineært}, gir $m>n$ likninger $m(t_i;\vb x)=y_i$, altså $A\vb x\approx\vb b$.
\begin{itemize}
  \item Rett linje $a+bt$: rad $i$ er $(1,\ t_i)$. Polynom $a+bt+ct^2$: rad $(1,t_i,t_i^2)$ (Vandermonde-matrise).
  \item Periodisk: $a+b\cos(2\pi t)+c\sin(2\pi t)$, som er lineær i $(a,b,c)$ selv om den ikke er lineær i $t$.
  \item Eksponentiell $ae^{bt}$: ta logaritmen, $\ln y=\ln a+bt$, som er lineær i $(\tilde a,b)$ med $a=\ee^{\tilde a}$. Da minimerer du avvik i $\ln y_i$, ikke i $y_i$, og svaret blir et annet enn med den ikke-lineære modellen.
  \item B-spline- og Bézier-tilpasning (F11) er også lineære i koeffisientene.
\end{itemize}

\begin{defn}[lineært minste kvadraters problem]
For $A\in\R^{m\times n}$ med $m>n$ og full rang $n$:
\[
\vb x^*=\argmin_{\vb x\in\R^n}\norm{A\vb x-\vb b}_2^2=\argmin\sum_{i=1}^m\bigl(m(t_i;\vb x)-y_i\bigr)^2 .
\]
\textbf{Tolkning:} $A\vb x^*=\sum_jx_j^*\vb a_j$ er det punktet i $\ran A$ (kolonnerommet) som ligger nærmest $\vb b$, altså den ortogonale projeksjonen. Residualet $A\vb x^*-\vb b$ står vinkelrett på alle kolonnene.
\end{defn}
\begin{teorem}[tre måter å løse det på]
\begin{enumerate}
  \item \textbf{Normallikningene:} $A^\T A\,\vb x=A^\T\vb b$. Entydig løsning når $\rank A=n$. Dette er betingelsen $\nabla\frac12\norm{A\vb x-\vb b}^2=A^\T(A\vb x-\vb b)=\vb 0$. Numerisk ulempe: $\operatorname{cond}(A^\T A)=\operatorname{cond}(A)^2$, for eksempel med Vandermonde-matriser.
  \item \textbf{Moore--Penrose:} $\vb x^*=A^\dagger\vb b$ (F1). Fungerer også uten full rang, og gir da minimum-norm-løsningen.
  \item \textbf{QR:} $A=QR$ med ortogonal $Q\in\R^{m\times m}$ og øvre triangulær $R$. La $\tilde R\in\R^{n\times n}$ være den øverste delen av $R$, sett $\vb d=Q^\T\vb b$ og løs $\tilde R\vb x=(d_1,\dots,d_n)^\T$ ved tilbakesubstitusjon. Kvadratfeilen er $\norm{A\vb x^*-\vb b}_2^2=d_{n+1}^2+\dots+d_m^2$. Det følger av at $\norm{A\vb x-\vb b}=\norm{Q^\T(A\vb x-\vb b)}=\norm{R\vb x-\vb d}$.
\end{enumerate}
\end{teorem}
\begin{eksempel}[øving 1, oppg.~3: linje gjennom $(0;1{,}25)$, $(1;3)$, $(2;4{,}8)$, $(3;6{,}9)$ (kontrollregnet)]
$A^\T A=\begin{psmallmatrix}4&6\\6&14\end{psmallmatrix}$ og $A^\T\vb b=\begin{psmallmatrix}15{,}95\\33{,}3\end{psmallmatrix}$ gir
$\vb x=\frac1{20}\begin{psmallmatrix}14&-6\\-6&4\end{psmallmatrix}\begin{psmallmatrix}15{,}95\\33{,}3\end{psmallmatrix}=\begin{psmallmatrix}1{,}175\\1{,}875\end{psmallmatrix}$. Linjen er $p(t)=1{,}175+1{,}875t$, med feil $\norm{A\vb x-\vb b}_2\approx0{,}184$.
\end{eksempel}
\begin{fallgruve}[overtilpasning]
Med mange parametre (polynom av høy grad) kan modellen treffe støyen i stedet for trenden. Et polynom av grad $m-1$ går gjennom alle $m$ punktene, men svinger mellom dem. I øving 1 skjer dette med kubisk modell.
\end{fallgruve}

\begin{oving}
\begin{itemize}[leftmargin=5.6em,labelwidth=5em,labelsep=0.6em,align=left]
\item[\lett] \href{https://github.com/gustavbv/Matte4C---eksamen/blob/main/oppgavesett/F12-minste-kvadrater/3-hele-oppgaven-i-pensum/009-TMA4115-2025-var-oppg2-oppgave.pdf}{TMA4115 vår 2025, oppg.~2} {\color{graa}(\href{https://github.com/gustavbv/Matte4C---eksamen/blob/main/oppgavesett/F12-minste-kvadrater/3-hele-oppgaven-i-pensum/009-TMA4115-2025-var-oppg2-svar.pdf}{LF})}: rett linje gjennom tre punkter
\item[\middels] \href{https://github.com/gustavbv/Matte4C---eksamen/blob/main/oppgavesett/F12-minste-kvadrater/1-TMA4432-ovinger-og-proveeksamen/002-TMA4432-oving-1-oppg3-oppgave.pdf}{Øving 1, oppg.~3} {\color{graa}(\href{https://github.com/gustavbv/Matte4C---eksamen/blob/main/oppgavesett/F12-minste-kvadrater/1-TMA4432-ovinger-og-proveeksamen/002-TMA4432-oving-1-oppg3-svar.pdf}{LF})}: lineær regresjon med Moore--Penrose
\item[\middels] \href{https://github.com/gustavbv/Matte4C---eksamen/blob/main/oppgavesett/F12-minste-kvadrater/3-hele-oppgaven-i-pensum/031-TMA4115-2018-kont-oppg9-oppgave.pdf}{TMA4115 kont.\ 2018, oppg.~9} {\color{graa}(\href{https://github.com/gustavbv/Matte4C---eksamen/blob/main/oppgavesett/F12-minste-kvadrater/3-hele-oppgaven-i-pensum/031-TMA4115-2018-kont-oppg9-svar.pdf}{LF})}: andregradspolynom tilpasset fem punkter
\item[\vanskelig] \href{https://github.com/gustavbv/Matte4C---eksamen/blob/main/oppgavesett/F12-minste-kvadrater/2-Bergmann-eksamener/006-TMA4215-2021-host-oppg1-oppgave.pdf}{TMA4215 høst 2021, oppg.~1} {\color{graa}(\href{https://github.com/gustavbv/Matte4C---eksamen/blob/main/oppgavesett/F12-minste-kvadrater/2-Bergmann-eksamener/006-TMA4215-2021-host-oppg1-svar.pdf}{LF})}: QR med Givens-rotasjoner og minste kvadrater
\item[\vanskelig] \href{https://github.com/gustavbv/Matte4C---eksamen/blob/main/oppgavesett/F12-minste-kvadrater/2-Bergmann-eksamener/005-TMA4215-2022-host-oppg2-oppgave.pdf}{TMA4215 høst 2022, oppg.~2} {\color{graa}(\href{https://github.com/gustavbv/Matte4C---eksamen/blob/main/oppgavesett/F12-minste-kvadrater/2-Bergmann-eksamener/005-TMA4215-2022-host-oppg2-svar.pdf}{LF})}: gradient descent sammenlignet med normallikningen
\end{itemize}
\end{oving}


\subsection{Ikke-lineære minste kvadrater}
\begin{defn}
Med et residual $\vb R(\vb x)=(r_1(\vb x),\dots,r_m(\vb x))^\T$ (for eksempel $r_i=m(t_i;\vb x)-y_i$):
\[
\argmin_{\vb x\in\R^n}E(\vb x),\qquad E(\vb x)=\tfrac12\norm{\vb R(\vb x)}_2^2=\tfrac12\sum_{i=1}^mr_i(\vb x)^2,
\]
\[
\nabla E(\vb x)=\vb R'(\vb x)^\T\vb R(\vb x),\qquad \Hess E(\vb x)=\vb R'(\vb x)^\T\vb R'(\vb x)+\sum_{i=1}^mr_i(\vb x)\Hess r_i(\vb x).
\]
\end{defn}
\begin{center}
\small
\begin{tabular}{@{}lll@{}}
\toprule
& lineær & ikke-lineær\\
\midrule
Residual $\vb R(\vb x)$ & $A\vb x-\vb b$ & vilkårlig\\
Jacobimatrise $\vb R'(\vb x)$ & $A$, konstant & avhenger av $\vb x$\\
$\nabla E=\vb R'^\T\vb R$ & $A^\T(A\vb x-\vb b)$, lineær & ikke-lineær\\
$\nabla E=\vb 0$ & normallikningene & ikke-lineært system\\
Løses med & ett lineært system & iterasjon\\
Eksempel & $a+bx+cx^2$ & $a\ee^{bx}$\\
\bottomrule
\end{tabular}
\end{center}

\begin{oppskrift}[Gauss--Newton]
Newton på $\nabla E=\vb 0$, men uten andrederiverte: $\Hess E\approx J^\T J$ med $J=\vb R'$. Fra $\vb x^{(0)}$, for $k=0,1,2,\dots$:
\begin{enumerate}
  \item Regn ut $\vb R(\vb x^{(k)})$ og $J^{(k)}=\vb R'(\vb x^{(k)})$.
  \item Løs $(J^{(k)})^\T J^{(k)}\Delta^{(k)}=-(J^{(k)})^\T\vb R(\vb x^{(k)})$ (med QR eller $A^\dagger$). Det er det samme som det lineære minste kvadraters problemet $\Delta^{(k)}=\argmin_\Delta\norm{J^{(k)}\Delta+\vb R(\vb x^{(k)})}_2^2$.
  \item $\vb x^{(k+1)}=\vb x^{(k)}+\Delta^{(k)}$. Stopp når $\norm{\vb R}$ eller $\norm{\Delta}$ er liten.
\end{enumerate}
\textbf{Svakheter:} metoden konvergerer bare lokalt og kan divergere. Lineariseringen $\vb R(\vb x+\Delta)\approx\vb R(\vb x)+J\Delta$ er dårlig for store $\Delta$, så $E$ kan til og med øke. $J^\T J$ kan mangle full rang, og da er løsningen ikke entydig.
\end{oppskrift}
\begin{oppskrift}[Levenberg--Marquardt (Gauss--Newton med tillitsområde)]
Straff store steg: $E_\lambda(\Delta)=\frac12\norm{J\Delta+\vb R}_2^2+\frac\lambda2\norm\Delta_2^2=\frac12\Bigl\|\begin{psmallmatrix}J\\ \sqrt\lambda I\end{psmallmatrix}\Delta-\begin{psmallmatrix}-\vb R\\ \vb 0\end{psmallmatrix}\Bigr\|_2^2$. Normallikningen blir
\[
\bigl(J^\T J+\lambda I\bigr)\Delta=-J^\T\vb R,
\]
og matrisen er inverterbar for $\lambda>0$. Stor $\lambda$ gir korte steg, omtrent gradient descent med steglengde $1/\lambda$. Liten $\lambda$ gir omtrent Gauss--Newton. $1/\lambda$ fungerer som en tillitsradius.\\
\textbf{Algoritmen:} gitt $\vb x^{(0)}$ og $\lambda_0>0$, løs for $\Delta^{(k)}$. Hvis $E(\vb x^{(k)}+\Delta^{(k)})<E(\vb x^{(k)})$: ta steget og sett $\lambda_{k+1}=0{,}8\lambda_k$. Ellers: bli stående og sett $\lambda_{k+1}=2\lambda_k$.
\end{oppskrift}

\begin{oving}
Oppgavebanken har ingen gamle eksamensoppgaver om Gauss--Newton eller Levenberg--Marquardt. Bruk denne til å øve:\\
\textbf{Oppgave.} Modellen $m(t;a,b)=a\ee^{bt}$ skal tilpasses $(t_i,y_i)=(0;2)$, $(1;3)$, $(2;5)$. (a) Skriv opp $\vb R$ og $J$. (b) Ta ett Gauss--Newton-steg fra $(a,b)=(2;0{,}5)$. (c) Hva endres i (b) med Levenberg--Marquardt og $\lambda=1$?\\
\textbf{Fasit (kontrollregnet):} $r_i=a\ee^{bt_i}-y_i$, og rad $i$ i $J$ er $(\ee^{bt_i},\ at_i\ee^{bt_i})$. I startpunktet er $\vb R\approx(0;\,0{,}2974;\,0{,}4366)$ og $J\approx\begin{psmallmatrix}1&0\\1{,}6487&3{,}2974\\2{,}7183&10{,}873\end{psmallmatrix}$. Gauss--Newton gir $\Delta\approx(-0{,}0768;\,-0{,}0235)$, altså $(a,b)\approx(1{,}923;\,0{,}477)$, og $E$ synker fra $0{,}140$ til $0{,}0077$. Med $\lambda=1$ blir $\Delta\approx(-0{,}0507;\,-0{,}0304)$, et kortere steg i en litt annen retning ($E\approx0{,}0083$).
\end{oving}
